
\documentclass[10pt]{article}
\usepackage{graphicx}
\usepackage{amsmath}
\usepackage{amssymb}
\usepackage{amsthm}
\usepackage{amsthm} 

\newtheorem{hypothesis}{Hypothesis} 
\usepackage{pdflscape}
\usepackage{afterpage}
\usepackage{rotating}
\usepackage{color}
\usepackage{placeins}
%\usepackage{paralist}
\usepackage{graphicx}
%\usepackage[retainorgcmds]{IEEEtrantools}
%\usepackage{hyperref}
\usepackage{comment} 
\usepackage{caption}
\usepackage{floatrow}
\floatsetup[table]{capposition=top}
\usepackage{subcaption}
\usepackage{booktabs}
\usepackage{natbib}
 \usepackage{tikz}
 \usepackage{changes}
\usetikzlibrary{decorations.pathreplacing}
\usepackage{longtable}
\usepackage[flushleft]{threeparttable}
\usepackage{framed}
%\usepackage{soul}
\usepackage{subfiles}
\bibpunct{(}{)}{;}{a}{}{,}
\graphicspath{{./plot/}}
\usepackage{array}
\newcolumntype{H}{>{\setbox0=\hbox\bgroup}c<{\egroup}@{}}

\makeatletter
\newcommand*{\textoverline}[1]{$\overline{\hbox{#1}}\m@th$}
\makeatother

%%% PAGE DIMENSIONS
 %%% PAGE DIMENSIONS
 \usepackage[margin=1.5in]{geometry}
 \geometry{a4paper}
\usepackage{setspace}
\doublespacing

\title{Culture, Coordination, and Team Performance}
\author{}
\date{\today}

\begin{document}


\maketitle
\thispagestyle{empty}

\begin{abstract}

xxx
\end{abstract}
\section{Introduction}
Organizations increasingly rely on culturally diverse teams to solve complex tasks and compete in global markets. While diversity expands the range of knowledge, perspectives, and expertise available within teams, it may also create communication barriers and increase coordination costs. Consequently, whether cultural diversity ultimately enhances or hinders team performance remains a central question in management and economics. Economic theory emphasizes the trade-off between the informational benefits of diversity and the communication frictions arising from heterogeneous backgrounds (Lazear, 1999), while organizational research similarly concludes that multicultural teams simultaneously benefit from broader informational resources but face greater communication difficulties and social integration challenges (Stahl et al., 2010). Organizational research further suggests that effective coordination depends not only on communication itself but also on the emergence of shared mental models—compatible understandings of the task, the environment, and teammates’ likely actions that allow individuals to anticipate one another without constant explicit communication (Cannon-Bowers, Salas and Converse, 1993; Klimoski and Mohammed, 1994; Mathieu et al., 2000). Cultural similarity may facilitate the development of such shared understandings by reducing communication frictions and providing teammates with more compatible expectations about one another’s behaviour. Yet despite extensive evidence on how culture influences communication and collaboration within organizations, comparatively little is known about whether cultural similarity ultimately translates into superior team performance.

In this paper, we study whether cultural similarity is associated with the performance of small, interdependent teams. We use professional men's Grand Slam doubles tennis as an empirical model of organizational coordination, where two players with distinct individual abilities must repeatedly coordinate their actions to produce a joint outcome.  Many organizations rely on dyadic or small-team collaboration, including project teams, surgical teams, military units, emergency response teams, and entrepreneurial partnerships. Although these settings differ substantially in purpose and context, they share a common organizational challenge: team members must repeatedly coordinate their actions under time constraints while adapting to rapidly changing environments. Doubles tennis captures these fundamental features in a particularly transparent way. Every point requires two teammates to coordinate positioning, movement, tactical decisions, and responses to their opponents, while success depends not only on individual ability but also on how effectively partners anticipate and complement each other's actions. At the same time, performance is measured objectively, tasks are standardized across teams, and teammates operate under strong incentives to maximize joint success. These characteristics make doubles tennis an informative empirical model of organizational coordination in environments where communication and mutual understanding are essential to collective performance.

\added{rewrite} We collect data from all men's  doubles matches in Grand Slams between 2018 and 2025. We measure cultural similarity along three complementary dimensions: shared nationality, shared official language, and linguistic proximity. These measures capture different aspects of common communication and cultural background. We estimate whether culturally similar teams outperform culturally dissimilar teams in completed matches as well as in tiebreaks where outcomes are determined over a very small number of points, limiting the opportunity for coordination advantages to accumulate  as a boundary condition. 

 We find that culturally similar teams are consistently more likely to win Grand Slam matches, with the strongest and most precisely estimated associations observed for the language-based measures of similarity. However, these advantages disappear in tiebreaks. Rather than contradicting the coordination hypothesis, this pattern suggests an important boundary condition. Tiebreaks represent condensed episodes in which the competing teams have already demonstrated closely matched performance and relatively few points determine the outcome. Under these conditions, coordination advantages have less opportunity to accumulate and short-run variation plays a greater role in determining the winner. Finally, we find that the performance association of shared language and linguistic proximity becomes significantly stronger as teammates accumulate professional experience, consistent with communication advantages developing gradually through repeated interaction rather than arising automatically from common cultural backgrounds.

Our study relates to four strands of literature. First, it contributes to the literature on cultural diversity and organizational performance. Economic models emphasize the trade-off between diversity's informational benefits and its communication costs (Lazear, 1999), while organizational research documents that cultural diversity simultaneously enhances creativity and knowledge variety but also creates coordination challenges and interpersonal conflict (Stahl et al., 2010; Harrison and Klein, 2007). Although numerous studies examine workforce diversity and firm outcomes, relatively few identify the role of cultural similarity within small teams performing highly interdependent tasks.

Second, our paper contributes to the literature on language, communication, and organizational coordination. Language has been recognized as a fundamental organizational resource that facilitates information exchange, knowledge integration, and collective decision-making (Wernerfelt, 2004). In international business and organizational research, language differences have been shown to impede knowledge transfer and collaboration across multinational organizations (Marschan-Piekkari, Welch and Welch, 1999; Feely and Harzing, 2003). Related evidence from economics shows that linguistic similarity reduces communication frictions and facilitates economic interaction, including trade, migration, and cross-border collaboration (e.g., Melitz and Toubal, 2014). Our analysis extends this literature by examining whether shared language is associated with performance in teams whose success depends on continuous real-time coordination rather than one-off communication.

Third, we relate to the literature on shared mental models and team coordination. Organizational researchers argue that high-performing teams develop common understandings of tasks and teammates that enable implicit coordination, particularly in environments characterized by repeated interaction and rapidly changing conditions (Cannon-Bowers et al., 1993; Klimoski and Mohammed, 1994; Mathieu et al., 2000). While this literature has primarily relied on laboratory experiments and organizational settings, our study provides evidence from elite professional teams performing under naturally occurring competitive conditions, where coordination quality is reflected directly in objective performance outcomes.

Finally, our paper builds on the emerging literature examining cultural homophily in professional sports. Most closely related is the recent Management Science study by Békés and Ottaviano (2025), which shows that professional football players exhibit systematic cultural homophily in passing behaviour, indicating that cultural similarity shapes patterns of collaboration even within highly professional organizations. Whereas that study focuses on who collaborates with whom, we examine whether cultural similarity is associated with how successfully teams perform once formed. By distinguishing collaboration from performance, we address a complementary question about the organizational consequences of cultural homophily.

This paper makes three contributions. First, it extends the emerging literature on cultural homophily by moving beyond collaboration patterns to examine objective team performance. Second, it contributes to the literature on communication and team coordination by showing that language-related dimensions of cultural similarity exhibit stronger and more persistent associations with performance than shared nationality alone, and that these associations become more pronounced as teams accumulate professional experience. Third, it identifies an important boundary condition for the value of cultural similarity. Rather than providing a universal performance advantage, cultural similarity appears to matter most in environments where teammates have repeated opportunities to coordinate over the course of a task and where communication advantages can accumulate over time, but much less in short, closely contested episodes such as tiebreaks, where outcomes are inherently more susceptible to point-level variation. More broadly, our findings suggest that cultural similarity is best understood not simply as a source of social affinity, but as a communication resource whose value depends on the opportunities available for coordination to influence performance.
\section{Data}
\subsection{Empirical Setting and Sample}
Our empirical setting consists of men's doubles matches from the four Grand Slam tournaments, the Australian Open, Roland Garros, Wimbledon, and the US Open, between 2018 and 2025. The dataset contains 1,876 completed Grand Slam matches.\footnote{From the initial dataset with 1,933  matches we exclude retirements and walkovers because these outcomes do not reflect normal competitive performance. We further remove matches with incomplete doubles ranking information that cannot be reliably recovered, resulting in a final estimation sample of 1,876 matches.} Since each match contributes one winning and one losing team,  the main analysis contains 3,752 units of team-match observations.

Player information is collected by combining official tournament records with ATP player information and supplementary public sources. Nationality and language information are fully harmonized through a multi-stage validation procedure that resolves inconsistencies arising from spelling variants, accented characters, historical nationality changes, and special cases such as Monaco. Only observations for which all four players have complete ranking and cultural information are retained.

\textcolor{red}{Table 1: Observation Breakdown by Grand Slams and Year and descriptive stats with average age, experience, ranking at team level}
\subsection{Measuring Cultural Similarity}
We examine three complementary measures of cultural similarity that capture different dimensions of communication and shared background.

The first measure, Same Nationality, equals one if both teammates represent the same country and zero otherwise. The second measure, Same Language, equals one if both teammates share an official national language. The third measure, Language Proximity, is a continuous measure ranging from zero to one that captures the linguistic similarity between teammates' native languages, with larger values indicating greater linguistic proximity.

These measures reflect increasingly broad notions of cultural similarity. Shared nationality captures common institutional and cultural background. Shared language more directly captures communication potential. Linguistic proximity recognizes that communication costs may also be lower between speakers of closely related languages even when they do not share the same official language.

\textcolor{red}{Table 2: Shares of three Cultural Similarity Measures}
Approximately X\% of teams consist of players sharing the same nationality,  X\% share an official language, and X\% share linguistic promixity.
\section{Theory and Hypotheses}
\subsection{Cultural Similarity and Team Coordination}
A central concept in the organizational literature is the idea of shared mental models. Teams develop shared representations of tasks, roles, priorities, and appropriate responses to different situations. When team members hold sufficiently similar understandings of these dimensions, they can anticipate one another's actions and coordinate without having to communicate every decision explicitly. Related work on implicit coordination emphasizes that effective teams can rely on shared expectations when responding to recurring situations. Transactive memory systems provide a complementary perspective: through repeated interaction, team members develop an understanding of what their teammates know and how knowledge and responsibilities are distributed within the team.

These perspectives suggest that coordination depends partly on reducing the cognitive and communication costs of joint action. Culture and language can affect these costs. Teammates who share a language may communicate intentions more easily and interpret verbal and non-verbal cues more consistently. Shared nationality may capture additional dimensions of cultural familiarity, including norms, expectations, and common social experiences. More generally, linguistic proximity may reduce communication barriers even when teammates do not share a nationality or an official language.

The implication is not that cultural similarity directly makes individual players better. Rather, cultural similarity may make it easier for teammates to combine their individual abilities into effective joint performance. This distinction is central to our analysis. We are interested in the coordination return to cultural similarity, rather than in differences in individual ability associated with nationality, language, or culture. We therefore propose:
\begin{hypothesis}
 Teams with greater cultural similarity are more like to win matches  than culturally dissimilar teams.
\end{hypothesis}
\subsection{Boundary Condition:  Tiebreaks}
The coordination benefits associated with cultural similarity should not necessarily be equally observable in every part of a match. Coordination is a process that can develop through repeated interaction: teammates communicate, adjust their behavior, learn one another's tendencies, and build shared expectations. A full match provides many opportunities for these processes to operate.

Tiebreaks provide a particularly informative boundary condition. \added{Add a short definition of tiebreaks here.} First, they are short relative to the match as a whole, leaving less time for a coordination advantage to accumulate. Second, a tiebreak is reached only when the two teams have performed sufficiently similarly to reach a closely contested stage of the set. Its outcome is therefore determined in a setting where relatively small differences in performance can be decisive and where point-level randomness is likely to play a larger role.

The communication environment also differs. During a set, teammates have repeated opportunities to communicate during changeovers, allowing them to exchange information, adjust tactics, and coordinate subsequent play. A tiebreak substantially compresses these opportunities. Once the tiebreak begins, teammates have very limited scope for formal discussion before the decisive points are played, with the principal changeover occurring only after the score reaches 6–6. Thus, a tiebreak provides not only a shorter horizon over which coordination can operate but also fewer opportunities to communicate and update a shared understanding during that period.

These features make the tiebreak a useful test of the coordination interpretation. If cultural similarity improves team performance partly by facilitating communication, shared expectations, and coordinated action, its advantage should be less apparent in a short, closely contested sequence in which there are fewer opportunities for communication and greater scope for point-level randomness. This leads to our second hypothesis:

\begin{hypothesis}
The performance advantage associated with cultural similarity is weaker in tiebreaks than in overall match outcomes.
\end{hypothesis}

\section{Empirical Analysis}

\subsection{Baseline Performance Effects}
We begin by testing Hypothesis 1 by examining whether culturally similar teams outperform culturally dissimilar teams in Grand Slam doubles. Our baseline specification estimates the association between cultural similarity of the two players within a team and the probability of winning a match. Specifically, we estimate
\begin{align}
  \operatorname{Pr}(W_{ijt} = 1) = f(\beta C_{ijt} + \gamma X_{ijt} + \delta_{ty} + \lambda_r) , 
\end{align}


where $W_{ijt}$ equals one if team $i$ wins match $j$ in tournament-year $t$, $C$ denotes one of our three measures of cultural similarity: same nationality, same language, and language proximity. %$X$ contains team and opponent characteristics which we elaborate further below, $\delta_{ty}$ denotes tournament-by-year fixed effects, and $\lambda_r$ denotes round fixed effects.


The variables in X  represent observable team and opponent qualities. Team quality is measured using the average ATP doubles ranking of the two teammates together with the average ranking of the opposing team. Lower ranking values indicate stronger players. We additionally include an indicator for whether either teammate has previously been ranked inside the ATP Top 100 in singles, capturing exceptional individual talent that may not be reflected fully in doubles rankings. Professional experience is measured as the average number of years since the two teammates turned professional. To allow for diminishing returns, we include both the demeaned experience measure and its squared term. Demeaning facilitates interpretation of interaction models while leaving the fitted values unchanged.

Throughout the analysis we include tournament-by-year fixed effects $\delta_{ty}$ and round fixed effects $\lambda_r$ to absorb common variation across tournaments and stages of competition. Standard errors are clustered at the match level.



\textcolor{red}{Regression Table: Match win}

Table X report average marginal effects (AMEs) estmated in logit models. which can be interpreted as percentage-point changes in the probability of winning associated with each explanatory variable. 

Across all three specifications, culturally similar teams are significantly more likely to win Grand Slam matches. The estimated effects range from 3.8 to 4.7 percentage points and are strongest when cultural similarity is measured through shared language rather than shared nationality. Team quality remains the dominant predictor of success, as expected: stronger doubles rankings significantly increase the probability of winning, while more experienced teams also perform better, although the experience-performance relationship is concave.

These findings provide consistent evidence that cultural similarity is positively associated with team performance after accounting for observable differences in team quality, opponent strength, experience, and tournament characteristics.

\subsection{A Boundary Condition: Tiebreak Performance} 
%Tiebreaks provide a particularly demanding boundary test because they are both short and conditional on the two teams having reached a closely contested state. The limited number of points leaves less scope for coordination advantages to accumulate, while the outcome of a short sequence is more sensitive to point-level variation.


%\begin{itemize}
%\item Compression: there are relatively few points over which coordination advantages can accumulate.
%\item Selection into a close contest: the tiebreak is reached when the two teams have been sufficiently evenly matched in that part of the match.
%\item Greater outcome variance: A short sequence of points is more susceptible to individual point-level events, so even if cultural similarity improves coordination, that advantage may be harder to detect in the winner of one tiebreak.
%\end{itemize}

To test Hypothesis 2, we estimate the same empirical specification using individual tiebreaks as the unit of observation.

\textcolor{red}{Regression Table: Tiebreaks win}

In contrast to the baseline analysis, none of the three measures of cultural similarity is significantly associated with the probability of winning a tiebreak. The estimated coefficients are economically small and statistically indistinguishable from zero. Importantly, the ranking controls remain highly significant, indicating that the specification continues to capture underlying differences in team quality.

Taken together, these findings suggest that the performance advantage associated with cultural similarity is not universal. Instead, it appears to emerge primarily over the course of an entire match, where teammates repeatedly coordinate across many points, rather than within short, closely contested episodes where relatively few points determine the outcome.
%A coordination advantage need not translate into a detectable advantage when the outcome is determined over a very short sequence.



\subsection{Robustness check: Testing Tennis-Specific Explanations} 


One concern is that the observed relationship reflects strategic characteristics unique to particular playing surfaces rather than general coordination processes. Grass courts reward shorter rallies and faster exchanges, whereas clay courts involve longer rallies and different tactical patterns. 
If the observed association were driven by a particular tactical or environmental feature of doubles tennis, we might expect the return to cultural similarity to vary across surfaces, which differ substantially in their playing characteristics. We therefore test whether the association differs between grass and other surfaces and between clay and other surfaces.


We interact each cultural similarity measure with surface. 


The coordination advantage associated with cultural similarity is remarkably stable across playing environments. The result is not driven by any specific physical environment of the tournament.

\section{Mechanisms}
\subsection{Coordination Through Shared Experience} 


On the one hand, experienced professionals have developed richer tactical knowledge and communication routines, allowing them to exploit shared language more effectively. 

\begin{itemize}

\item Does cultural similarity become more valuable as teams accumulate experience? 
\end{itemize}
We test this mechanism by augmenting the baseline specification by interacting each measure of cultural similarity with average years since turning professional.

We estimate the equation 
\begin{align}
\operatorname{Pr}(\text{Win}) = f(\beta_1 C + \beta_2 \text{AvgTeamExp} + \beta_3 (C \times \text{AvgTeamExp}) + \gamma X + \delta_{ty} + \lambda_r)
\end{align}
\subsection{Knowledge Transfer with Experience Gaps} 

On the other hand, shared language becomes especially valuable when teammates differ substantially in expertise because communication and knowledge transfer become more demanding. Large experience gaps imply asymmetric knowledge, asymmetric decision authority, coaching/mentoring,
 and more need to transfer tacit knowledge. That is exactly where shared language may matter most.
\begin{itemize}
\item Does cultural similarity matter particularly when teammates differ substantially in experience?
\end{itemize}
The estimating equation becomes
\begin{align}
\operatorname{Pr}(\text{Win}) = f(\beta_1 C + \beta_2 \text{TeamExpGap} + \beta_3 (C \times \text{TeamExpGap}) + \gamma X + \delta_{ty} + \lambda_r)
\end{align}
%There may be greater uncertainty about: how the less experienced player behaves; tactical preferences; positioning; risk-taking; communication; expectations about the partner. Shared language could potentially reduce those coordination frictions.
\subsection{Alternative Explanation: Individualism vs Collectiveness}
Is the value of cultural similarity greater when the underlying cultural environment places more emphasis on collectivism? We find no evidence.

 

\section{Descriptive Evidence on Team Formation}
\subsection{Partner Selection}
\subsection{Evidence from the Olympics}

If cultural similarity is associated with better performance, an important question is whether players appear to sort toward culturally similar partners when the incentives to coordinate are particularly salient. We therefore examine team composition around the Tokyo and Paris Olympic cycles.  

\textcolor{red}{Table X reports the composition of teams across Olympic preparation cycles.} Despite widespread speculation that Olympic eligibility may encourage players to form same-nationality partnerships, the share of same-nationality teams remains remarkably stable during the Tokyo Olympic cycle. The Tokyo preparation period coincides with the covid period. During that period, international travel and international interaction were unusually constrained.

If anything, this creates a plausible setting in which same-nationality partnering could have been especially attractive: easier travel; fewer international logistical complications; greater opportunities to train together; potentially greater value from existing social/professional networks. Even in a period when practical constraints may have made same-nationality partnering particularly convenient, we do not observe an increase in same-nationality team formation.


During the Paris cycle, same-nationality and same-language partnerships become modestly more common, although the increase occurs gradually rather than as a discrete shift during the Olympic preparation period. These patterns do not provide evidence of systematic optimization.


\section{Conclusion}
\end{document}
